%% =====================================================================
%%  T-14-07  The Label
%%  -------------------------------------------------------------------
%%  One idea per file. Core is mandatory. Slots in canonical order.
%%  Max 700 words. No layout commands. No filler. New source material
%%  for this entry goes in sources/<BOOK>/T-14-07.tex -- never by
%%  rewriting this file (docs/RULES.md R0.1).
%% =====================================================================
%% @kind: topic
%% @id: T-14-07
%% @title: The Label
%% @subtitle: Name the emotion and it loses the steering wheel
%% @chapter: c14
%% @order: 70
%% @status: stable
%% @sources: B5
%% @updated: 2026-09-19

\topic{T-14-07}{The Label}{Name the emotion and it loses the steering wheel}

\begin{core}
State the counterpart's emotion or stake as an observation --- some variant of \emph{it seems like} or \emph{it sounds like} --- and never as a question about them. Naming an emotion moves the counterpart from feeling it to confirming or correcting your description; either way they are now processing it, not obeying it.
\end{core}
\begin{mechanism}
An unlabelled emotion runs the person; a labelled one becomes an object of joint attention. The label works because it requires no confession: the counterpart can accept it silently, deny it, or refine it, and every path yields information while the framing stays theirs. Used on fear or anger, it defuses --- the counterpart hears their worst case said aloud without catastrophe following. Used on positives, it reinforces them. The pre-emptive form, the accusation audit, lists every negative the counterpart could think about you before they do: each accusation, spoken first by you, arrives with no defence attached and therefore disarms as mere fact.
\end{mechanism}
\begin{conditions}
Strongest under visible emotion --- fear, anger, resentment, suspicion --- and at the start of high-stakes contact, where the audit variant belongs. Weakest against flat, professional counterparts who feel nothing in play, where the label is just noise; and in cultures or relationships where naming feelings reads as intrusion, it costs more than it buys. Mislabelling is the operational risk: a wrong label tells the counterpart you do not understand them at all.
\end{conditions}
\begin{application}
  \item Infer the emotion from what they do, not what they say: hostages in a siege worry about dying and prison whether or not they mention it.
  \item Open with the two least threatening words --- \emph{it seems like} --- and name the stake: the fear, the constraint, the position you suspect they hold.
  \item Stop. Let them confirm, correct, or deny. Every response is intelligence.
  \item Re-label as the emotion moves. A label is a probe, not a verdict.
  \item For the audit form, open with the worst they could say about you --- \emph{you're going to think I'm unreasonable, and you have no reason to trust me} --- and pause. Let them waive it.
\end{application}
\begin{feedback}
  \item Landing: they answer the label's content --- \emph{exactly}, or \emph{no, the real issue is...} --- and the temperature drops.
  \item Landing on an audit: the counterpart dismisses the accusation and often volunteers reassurance, which reverses the frame.
  \item Failing: flat agreement with no elaboration, or a fast \emph{don't tell me how I feel}.
  \item Failing: they begin labelling you back --- the mirror has been raised.
\end{feedback}
\begin{failure}
A wrong label tells the counterpart you were performing understanding, and performance is worse than silence: it converts sympathy into suspicion. Labels used mechanically --- sprinkled to steer --- read as therapy-script handling, and the counterpart feels managed rather than met. The audit overdone sounds like a con man's rehearsed humility; two or three honest negatives outperform six polished ones.
\end{failure}
\begin{countermeasures}
  \item Notice the form: \emph{it seems like} almost never precedes a genuine request for your view; it precedes steering. Treat the label as a probe and answer with facts, not feelings.
  \item Refuse the frame politely: name what they are doing --- \emph{you're labelling me} --- and the tool is spent.
  \item When someone pre-lists your criticisms of them, ask yourself why they needed you not to hold them; the audit itself is the information.
  \item Keep your own stakes spoken plainly. Labels feed on the unsaid.
\end{countermeasures}

\sources{B5}
\seesources
\seealso{T-22-02, T-19-01, T-23-02}
